\begin{textsample}{NEG Dim 4 – vem\_ed\_23 – Score 1.00 – vem\_ed\_23/t117.txt}  \label{ex:f4_neg_087}
Leitores Aos Osvaldo Succi Jr.
Coordenador PCIs Com o amadurecimento dos Projetos Colaborativos Internacionais realizados nas Fatecs sob a coordenação da equipe de PCIs / Cesu, aumentam produções acadêmicas relatando reflexões e práticas sobre essas experiências.
Este número 23 de VEm destaca três exemplos: Participação da equipe dos PCIs / Cesu na Conferência Internacional Faubai 2024 com apresentação de três trabalhos: dois sobre temas emergentes – uso de técnicas de Big Data e Inteligência Artificial na avaliação e no planejamento de projetos COIL (Collaborative Online International Learning – e um sobre o papel dos professores de línguas estrangeiras na internacionalização.
Regiane Souza Camargo Moreira, responsável pelos PCIs em espanhol, fez comunicação oral resultante do projeto conduzido em parceria com Mónica Paola Díaz Oliveros (Uniminuto / Colômbia) em evento promovido pela Consejeria de Educación de la Embajada de España en Brasil e pelo Colégio Miguel de Cervantes.
Pedro Rosa, professor da Fatec Tatuí, desenvolve PCI / Cesu desde 2021 com BridgeValley Community College (EUA).
Nesses três anos, aprimorou sua competência linguística no inglês e conseguiu publicar no Journal of Physics (2024).
Além disso, relatamos as visitas dos professores Rob Steel (agosto de 2023) e Cathy Lee (março de 2024), que concedeu entrevista à jornalista Patrícia Patrício em 9 de maio de 2024.
Boa leitura!

% matched lemmas: 
\end{textsample}
